\subsection{ZERO-DIMENSIONAL SPACES AND LUSIN SPACES}

DEFINITION 5. A topological space is said to be zero-dimensional if it is Hausdorff and if every point has a fundamental system of neighbourhoods which are both open and closed.

Every zero-dimensional space X is totally disconnected; for the component of a point x is contained in all the sets containing x which are both open and closed (Chapter I, § 11, no. 5), and the intersection of these sets is just \(\{x\}\) if \(\mathbf{X}\) is zero-dimensional.

Conversely, a locally compact totally disconnected space is zero-dimensional (Chapter II, § 4, no. 4, Proposition 6); but there exist totally disconnected metrizable spaces which are not zero-dimensional {[}Exercise 3 b){]}.

Every subspace of a zero-dimensional space is zero-dimensional, and topological sums and products of zero-dimensional spaces are zero-dimensional.

DEFINITION 6. A topological space X is a Lusin space if it is metrizable and if there exists a zero-dimensional Polish space P and a continuous bijective mapping of P onto X.

Every Lusin space is clearly a Souslin space.

PROPOSITION 12. A metrizable space is a Lusin space if and only if it is the image of a Polish space under a continuous bijective mapping.

The condition is clearly necessary; let us show that it is also sufficient. If f is a continuous bijection of a Lusin space X onto a metrizable space Y, it follows from Definition 6 that Y is a Lusin space. Hence we need only show that a Polish space is a Lusin space.

Notice first that, if X is a Lusin space, every closed (resp. open) subspace A of X is a Lusin space (cf.~no. 7, Theorem 3); for if f is a continuous bijection of a zero-dimensional Polish space P onto X, then \(\overline{f}^{1}(A)\) is closed (resp. open) in P and is therefore a zero-dimensional Polish subspace of P (no. 1, Propositions 1 and 2).

Every countable product of Lusin spaces is a Lusin space; this follows from Proposition 1 of no. 1 and the fact that every product of zero-dimensional spaces is zero-dimensional. Every countable intersection of Lusin subspaces of a Hausdorff topological space is a Lusin subspace; this follows from the preceding remarks and from Lemma 1 of no. 1. Furthermore:

LEMMA 2. If a metrizable space X is such that there exists a countable partition \((\mathbf{A}_{n})\) of X formed of Lusin subspaces, then X is a Lusin space.

For each n, let \(P_{n}\) be a zero-dimensional Polish space and let \(f_{n}\) be a continuous bijection of \(P_{n}\) onto \(A_{n}\) . If P is the topological sum of the \(P_{n}\) , then P is Polish (no. 1, Proposition 1) and zero-dimensional, and the mapping \(f: P \to X\) which agrees with \(f_{n}\) on \(P_{n}\) for each n is a continuous bijection of P onto X; hence the result.

Let us now show that the interval \(I = [0, i]\) of R is a Lusin space. Let J be the subspace of I consisting of all irrational numbers; then J is Polish (no. 1, Corollary to Proposition 3); also J is zero-dimensional, for if x is any point of J, the traces on J of intervals of R of the form {]}r, s{[}, where r and s are rational and r \textless{} x \textless{} s, form a fundamental system of open and closed neighbourhoods of x in J (because the traces on J of {]}r, s{[} and {[}r, s{]} are the same). Hence J is a Lusin subspace of I. Now J and the subspaces of I which consist each of a single rational point form a countable partition of I, and therefore I is a Lusin space by Lemma 2.

Let P be an arbitrary Polish space. By Corollary 1 to Theorem 1 (no. 1), P is homeomorphic to a subspace of the cube IN which is a countable intersection of open sets in IN. Since I is a Lusin space, the remarks at the beginning of the proof show that P is a Lusin space, and the proof of Proposition 12 is therefore complete.

